\documentclass[12pt,a4paper]{article}
\usepackage[UTF8]{ctex}
\usepackage{amsmath}
\usepackage{geometry}
\usepackage{abstract}
\usepackage{graphicx}
\usepackage{enumitem}
\usepackage{float}
\usepackage{fancyhdr}
\usepackage{listings}
\usepackage{xcolor}       % 定义和使用自定义颜色
\usepackage{booktabs}
\numberwithin{figure}{subsection}
\numberwithin{figure}{subsection}

\special{dvipdfmx:config z 0} %取消PDF压缩，加快速度，最终版本生成的时候最好把这句话注释掉

\geometry{a4paper,left=2.5cm,right=2.5cm,top=2.5cm,bottom=2.5cm}

% 页眉页脚设置
\pagestyle{fancy}        % 使用 fancyhdr 宏包的页眉页脚样式
\fancyhf{}               % 清除默认页眉页脚
\fancyfoot[C]{\thepage} % 设置页脚中间部分
\renewcommand{\headrulewidth}{0pt}

% 设置字体
\setCJKmainfont{SimSun}[BoldFont=SimHei,Mapping=latin2]
\setCJKsansfont{SimHei}
\setCJKmonofont{SimSun}

% 标题格式
\title{\zihao{3}\heiti 标题（基于xxx的研究）}
\author{}
\date{}

% 章节格式
\usepackage{titlesec}
\titleformat{\section}
  {\zihao{4}\heiti\centering}{\thesection}{0.5em}{}
\titleformat{\subsection}
  {\zihao{-4}\heiti\raggedright}{\thesubsection}{0.5em}{}
\titleformat{\subsubsection}
  {\zihao{-4}\heiti\raggedright}{\thesubsubsection}{0.5em}{}

%摘要格式
\renewcommand{\abstractnamefont}{\heiti\bfseries\fontsize{16}{18}\selectfont}
\renewcommand{\abstracttextfont}{\songti\fontsize{12pt}{18}}

%调整图片标题与下文距离
\setlength{\belowcaptionskip}{-1cm}

%设置代码块格式
\lstset{
    basicstyle=\ttfamily\footnotesize, % 设置代码块的基础字体
    keywordstyle=\color{blue},         % 设置关键字的颜色
    stringstyle=\color{red},           % 设置字符串的颜色
    commentstyle=\color{gray},         % 设置注释的颜色
    backgroundcolor=\color{lightgray!20}, % 设置代码块背景色
    frame=single,                      % 添加单边框
    breaklines=true,                   % 自动换行
    numbers=left,                      % 显示行号
    numberstyle=\tiny\color{gray},     % 设置行号样式
    captionpos=b,                      % 标题放在代码块下方
}

% 正文
\begin{document}

% 第一页
\maketitle
\thispagestyle{empty}
\vspace{-2.5cm}
\begin{abstract}
\vspace{-0.5cm}
\zihao{-4}\songti
第一段，主要的概括...... 

针对问题一，采用了什么方法，确立了什么模型，得到了什么结果。

针对问题二，采用了什么方法，确立了什么模型，得到了什么结果。

针对问题三，采用了什么方法，确立了什么模型，得到了什么结果。

本文的特点：1.内容。2.内容。
\end{abstract}

\vspace{0cm}
\noindent%取消
\textbf{关键词：} 关键词1，关键词2，关键词3
\newpage

%插入目录
\tableofcontents
\thispagestyle{empty}
\newpage
\setcounter{page}{1}
% 主要内容
\section{问题重述}
\subsection{问题背景}
问题背景内容.....

插入图片
\begin{figure}[ht]%here or top
  \centering
  \includegraphics[width = 0.55\textwidth]{figures/Miku_Space.jpg}
  \caption{Miku in Space}
  \label{img:Miku}    
\end{figure}
\subsection{问题提出}
围绕xx的问题，本文主要研究以下几个问题：

问题一：xxx， 以下几个小问题：
\begin{enumerate}[itemsep=-0.2em,topsep=0em]
  \item 第一项
  \item 第二项
  \item 第三项
\end{enumerate}

问题二：xxx，以下几个小问题：
\begin{enumerate}[itemsep=-0.2em,topsep=0em]
  \item 第一项
  \item 第二项
  \item 第三项
\end{enumerate}

\section{模型假设}
为了简化模型，本文提出了以下假设：
\begin{enumerate}[itemsep=-0.2em,topsep=0em]
  \item 第一项
  \item 第二项
  \item 第三项
\end{enumerate}

\section{符号说明}
\vspace{-2em}

\begin{table}[H]
  \centering%把表居中
  \begin{tabular}{ccc}
  \toprule%第一道横线
  符号&说明&单位\\
  \midrule%第二道横线 
  $F_e$ & 波浪鼓励力 & $N$ \\
  \bottomrule%第三道横线
  \end{tabular}
\end{table}

注：表中未出现的符号及重复的符号以出现处的解释为准
\section{问题一模型的建立与求解}
\subsection{问题分析}
问题一包含如下三个子问题：
\begin{enumerate}[itemsep=-0.2em,topsep=0em]
  \item 问题一
  \item 问题二
  \item 问题三
\end{enumerate}
\begin{figure}[H]%here or top
  \centering
  \includegraphics[width = 0.55\textwidth]{figures/Frieren.jpg}
  \caption{Frieren}
  \label{img:Frieren}    
\end{figure}

\subsection{模型建立}
建立大致模型
\subsubsection{具体建立模型}
具体内容
\subsubsection{模型总述}
总结模型并写下公式
\begin{equation}
  \int_a^b f(x) \, dx = F(b) - F(a)
\end{equation}
\subsection{模型求解}
利用公式求解
\begin{equation}
  E=mc^2
\end{equation}
求解结果+图片
\begin{figure}[H]%here or top
  \centering
  \includegraphics[width = 0.55\textwidth]{figures/Frieren_Night.jpg}
  \caption{FrierenNight}
  \label{img:Frieren_Night}    
\end{figure}

\subsection{方案分析}
一些分析，之后可以继续改进。

\section{问题二模型的建立与求解}
\subsection{问题分析}
问题二包含如下三个子问题：
\begin{enumerate}[itemsep=-0.2em,topsep=0em]
  \item 第一项
  \item 第二项
  \item 第三项
\end{enumerate}
\begin{figure}[H]%here or top
  \centering
  \includegraphics[width = 0.55\textwidth]{figures/Miku_Dance.jpg}
  \caption{MikuDance}
  \label{img:MikuDance}    
\end{figure}

\subsection{模型建立}
建立大致模型
\subsubsection{具体建立模型}
具体内容
\subsubsection{模型总述}
总结模型并写下公式
\begin{equation}
  \int_a^b f(x) \, dx = F(b) - F(a)
\end{equation}
\subsection{模型求解}
利用公式求解
\begin{equation}
  E=mc^2
\end{equation}
求解结果+图片
\begin{figure}[H]%here or top
  \centering
  \includegraphics[width = 0.55\textwidth]{figures/Makima.png}
  \caption{Makima}
  \label{img:Makima}    
\end{figure}

\subsection{方案分析}
一些分析，之后可以继续改进。
\section{模型总结与评价}
\subsection{模型优点}
\begin{enumerate}[itemsep=-0.2em,topsep=0em]
  \item 问题一
  \item 问题二
  \item 问题三
\end{enumerate}

\subsection{模型缺点}
\begin{enumerate}[itemsep=-0.2em,topsep=0em]
  \item 问题一
  \item 问题二
  \item 问题三
\end{enumerate}

\subsection{模型改进}
\begin{enumerate}[itemsep=-0.2em,topsep=0em]
  \item 问题一
  \item 问题二
  \item 问题三
\end{enumerate}

\newpage
% 参考文献
\begin{thebibliography}{99}
\zihao{-4}\songti
\bibitem{ref1} 作者，书名，出版地：出版社，出版年。
\bibitem{ref2} 作者，论文名，杂志名，卷期号：出版年，起止页码。
\bibitem{ref3} 作者，资源标题，网址，访问时间（年月日）。
\end{thebibliography}

\newpage
\section*{附录A 数据支撑}
利用表格进行数据支撑
\section*{附录B 代码附录}
\begin{lstlisting}[language=Matlab]
  f = [-4000 -3000]';
  A = [2 1;1 1;0 1];
  b = [10 8 7]';
  lb = [0 0]';
  ub = [+inf;+inf];
  [x,val] = linprog(f,A,b,[],[],lb,ub)
  \end{lstlisting}
\end{document}
